\subsection{Observed disagreement between domain classifiers}\label{sec:adult}
A separate comparison uses observed outputs from the adult-domain study of
\citet{sood2024holier}. Its public domain summary records \AdultVisits{} visits
to \AdultDomains{} domains in the June 2022 YouGov panel. The study used a
Piedomains rule that marks a domain when its predicted class is ``porn'' and
its reported class probability is at least 0.033. We compare that rule with
recorded ``porn'' categories from Bitdefender, ``sex'' from Forcepoint, and
``Pornography'' from alphaMountain. Each row retains only domains with a
recorded category from that provider, and visit shares use the visits in that
row's frame as their denominator. Missing provider categories are not negative
labels.

\begin{table}[htbp]
\centering\small
\setlength{\tabcolsep}{3.5pt}
\caption{Observed classifier disagreement and visit-weighted shares}
\label{tab:adult}
\begin{tabular}{lrrrrrr}
\toprule
Reference & Domains & Visits (\%) & Agree (\%) & Domain + (\%) & Visit + (\%) & Count range (\%) \\
\midrule
Bitdefender & 43,391 & 84.2 & 97.7 & 1.92/3.36 & 1.05/1.71 & [0.07, 82.11] \\
Forcepoint & 47,950 & 87.5 & 98.0 & 1.65/3.04 & 1.02/1.64 & [0.07, 80.61] \\
alphaMountain & 36,063 & 52.7 & 99.0 & 0.98/1.35 & 1.38/2.16 & [0.04, 68.07] \\
\bottomrule
\end{tabular}

\par\smallskip
\begin{minipage}{\linewidth}\footnotesize
Piedomains is the proxy in every row; the named provider supplies the reference
category. ``Visits'' is the share of all \AdultVisits{} visits represented by
domains with a recorded reference category. Agreement counts domains equally.
``Domain +'' and ``Visit +'' list Piedomains/provider percentages, first for
domains counted equally and then for visits within that row's frame. The
count range is the sharp range of the provider visit share if only the two
directional disagreement counts are kept, holding Piedomains labels and visit
weights fixed. The providers use
different category definitions; none is an adjudicated truth label.
A recorded category other than the named positive one counts as negative for
this comparison, including nonspecific categories such as ``misc.''
\end{minipage}
\end{table}
\FloatBarrier

On the \AdultFocalDomains{} domains categorized by Bitdefender, which account
for \AdultFocalCoverage{}\% of visits, the two classifications agree on
\AdultFocalAgreement{}\% of domains. Piedomains marks
\AdultFocalProxyDomains{}\% of domains positive, compared with
\AdultFocalReferenceDomains{}\% under Bitdefender's category. Their
visit-weighted shares are \AdultFocalProxy{}\% and
\AdultFocalReference{}\%, respectively. The discrepancy is visible before
weighting domains by visits, so this comparison does not isolate an effect of
visit concentration. There are \AdultFocalProxyOnly{} Piedomains-only and
\AdultFocalReferenceOnly{} Bitdefender-only positive domains.

With both classifications observed, the provider's visit share is known. If
the disagreement locations were discarded and only their directional counts
retained, its compatible range would be
[\AdultFocalLower{}, \AdultFocalUpper{}]\%. A single Piedomains-negative
domain carries \AdultFocalLargest{}\% of visits in this frame and can dominate
the count-only upper bound. Appendix~\ref{app:adulthead} shows how that range
changes when the observed label pairs of the most-visited domains are retained.
This isolates information discarded by a confusion table; it does not
establish that the extreme allocations are plausible classifier errors.

The comparison measures disagreement between classification systems whose
categories differ. It does not estimate either system's error against
verified content labels. The domain-level comparison does not provide
regression SEs for individual browsing outcomes. Applying the item-audit
correction would require compatible person--domain links, verified labels
under a declared content rule, and the label-selection probabilities.
The Piedomains score and provider categories alone do not establish the
conditional probabilities required in Section~\ref{sec:calibration}.
